% ED26-BM-RAT2021-15-1 C01--C05: Markovits2021 developmental treatment.
\chapter{The development of basic human rationality}
\label{sec:p19-rationality-15-1}

Developmental evidence does not support identifying rationality with a single
score that inevitably increases with age. A child can draw an appropriate
inference without explaining its necessity; an adolescent can identify a
fallacy in a neutral argument yet overlook it in a favored argument; an adult
can understand validity but fail on a demanding inference. Markovits's 2021
chapter organizes these differences through explicit and implicit rationality,
and compares explanations based on reflection, social interaction, biological
specialization, cognitive capacity and semantic knowledge. Its conclusion is
an interaction account, not an agreed universal developmental timetable
\cite[pp.~779--787]{Markovits2021DevelopmentRationality}.

\begin{hbdefinition}{def:p19-r151-explicit-implicit}{Explicit and implicit rationality in the selected chapter}
\emph{Explicit rationality} concerns reasons offered for an inference, or
explicit assessments of why one inference is more adequate than another.
\emph{Implicit rationality} concerns the adequacy of judgments or inferences
actually produced, whether or not the participant gives an account of their
adequacy. \emph{Metalogical understanding} makes features such as inferential
validity and logical necessity objects of reflection. These are different
constructs, not three interchangeable names for an observed correct response.

The chapter mostly evaluates inference against standard logical norms,
because that is the orientation of much of the developmental work it reviews.
It also discusses practical, contextual and social adequacy. A task must
therefore specify its criterion: deductive necessity given stipulated premises
is not the same as empirical truth, a plausible causal expectation, or success
in a social exchange \cite[p.~779]{Markovits2021DevelopmentRationality}.
\end{hbdefinition}

\section{Reflection, argumentation and the use of competence}
\label{sec:p19-r151-explicit}

Moshman's account distinguishes making an inference from reasoning under an
understanding of its goals and standards. Reflection on particular inferences
can make their common structure explicit; increasingly general epistemic
understanding can continue into adulthood without a necessary endpoint.
The developmental claim concerns the construction and use of metalogical
concepts, not the first appearance of every valid inference. Correct modus
ponens is compatible with an inability to explain the difference between truth
and validity. Conversely, an explicit concept does not ensure its successful
application in every situation. Markovits identifies the relation between
developing explicit understanding and actual inference as an unresolved
empirical dependency. The adult confidence studies he cites do not provide a
longitudinal developmental test of that dependency
\cite[pp.~779--780]{Markovits2021DevelopmentRationality}.

Argumentation adds a social route: inferences and justifications are organized
to defend a position or challenge another. The social-constructivist account
suggests that encountering conflicting positions can prompt reflection on
relative inferential strength. The chapter reviews conservation interactions,
collaborative reasoning and Kuhn's argumentation interventions as evidence
relevant to this route. Those studies have different outcomes: a conservation
judgment, a selection-task answer, a dialogue strategy, and an individual essay
are not one measure of general rationality. A claim of improvement must name
which changed, under which interaction, and whether transfer was tested
\cite[pp.~780--781]{Markovits2021DevelopmentRationality}.
Markovits locates the contrast partly in the locus of development: individual
reflection on inferential structure versus socially produced confrontation of
positions. Both can advance explicit understanding without completely
determining subsequent judgments.

The argumentative account attributed to Mercier and Sperber gives a different
role to social interaction: argument production can serve persuasion and an
individual's position rather than impartial appraisal. Markovits contrasts
this with the expectation that confrontation necessarily leads toward better
explicit standards. His adolescent belief-bias example, attributed to
Klaczynski (2001), describes stronger detection of defects in relatively neutral
arguments than in arguments favoring emotionally held beliefs. It illustrates
context-sensitive deployment, not a demonstration that all adolescents lack
the relevant skill or that persuasion and truth are always opposed. These are
the chapter's theoretical and qualitative study comparisons; neither account
is established solely by the existence of a biased response
\cite[pp.~780--781]{Markovits2021DevelopmentRationality}.

% C05 continuation for ED26-BM-RAT2021-15-1. New reader material only.
% Root-owned placement: after the core's explicit-rationality/argumentation
% discussion. Existing core, prior evidence fragment and P11 record unchanged.
\subsection{Collaborative selection is not yet developmental transfer}
\label{sec:p19-r151-collaborative-interface}

Moshman and Geil's 1998 original experiment already has its full methodological
and numerical record in \cref{sec:p11-ri-individual-social}, with response-unit
conventions at \cref{def:p11-ri-collaborative-statistics}. It is the same
experiment cited in the developmental chapter, not a second independent
demonstration. Course-based undergraduates completed an abstract four-card
task individually or in consensus-oriented groups of five or six, with
written explanations. The original source gives 143 participants but compares
32 individual outputs with 20 group outputs. Its pooled 15/20 correct group
solutions and 3/32 correct individual solutions are not an age comparison
or 143 independent group trials
\cite[pp.~233--236]{MoshmanGeil1998CollaborativeReasoning}.

The developmental relevance lies in a possible mechanism: participants can
combine partial insights into a justified argument. In the two illustrated
groups, an argument for checking the odd-number card was combined with
an argument against checking the even-number card; neither group initially
contained a member choosing the complete correct pattern. Such construction
is different from counting votes or copying an initially complete answer.
Those illustrations were purposively selected from three all-initially-
incorrect groups; the third videotape was lost. The paper supplies no
formal process-coding reliability or isolated mediation estimate
\cite[pp.~237--245]{MoshmanGeil1998CollaborativeReasoning}.

The original discussion nevertheless does not establish a general childhood
developmental effect. Participant ages and a child cohort are not supplied.
Responses are immediate, on the same task, following discussion; no delayed
or novel-task test demonstrates lasting individual transfer. Group discussions
took approximately two-and-a-half times the individual solving time, without
a time-yoked control. Eliminating simple adoption of the initial majority
answer does not eliminate every influence of confidence, status or later
consensus. These limits accompany the existing original-study record rather
than being replaced by a generic claim that collaboration improves rationality.

At p.~246 the authors explicitly decline to claim durable, generalizable
competence change from this brief interaction. Their prediction that groups
of children younger than eleven or twelve might fare no better than individual
children, if the task requires a particular metalogical competence, is a
conditional prediction, not an observed child result. Repeated dialogic
learning is a proposed developmental route. The two adolescent interventions
below test related but different outcomes; they cannot retroactively turn
the undergraduate selection experiment into childhood longitudinal evidence
\cite[pp.~242--246]{MoshmanGeil1998CollaborativeReasoning}.

\subsection{Kuhn and Udell's 2003 intervention: discourse and its individual product}
\label{sec:p19-r151-kuhn-udell}

Kuhn and Udell distinguish an individual argument supporting a claim from
argumentative discourse responding to an interlocutor. Their intervention
targets a shift from articulating one's own position to engaging the
opponent's reasons, including critique and rebuttal. The assessment criterion
is a source-defined discourse/argument rubric, not an answer key identifying
the ethically correct position on capital punishment. More coded
counterarguments do not by themselves prove greater logical validity
\cite[pp.~1245--1248]{KuhnUdell2003Development}.

The analyzed sample comprised 34 eighth graders from two low-performing
New York City public middle schools. The article describes its target group
as academically at-risk thirteen-to-fourteen-year-olds but gives no individual
age distribution. It reports 14 Hispanic, 19 African-American and one
Ethiopian participant; 11 male and 23 female. These are the reported
sampling descriptors, not diagnoses or population-wide competence judgments.
There were 21 completers in the experimental condition and 13 in comparison.
Five additional starters, three experimental and two comparison, did not
complete because of chronic absence or suspension. Thus 39 started and
34 completed; outcome-specific samples are smaller in some analyses.

The authors report random assignment to conditions without specifying its
algorithm. Teams were then formed from pro/con opinions. The experimental
teams were represented at both schools; the comparison teams were at one
school. Team sizes were four to eight, with typical attendance three to
six. Recruitment/response rates, consent details and native-language
eligibility are not reported. English task formulations are supplied.
School, team and discussion-partner dependence remain relevant to inference
\cite[pp.~1247--1250]{KuhnUdell2003Development}.

The full condition received sixteen sessions scheduled twice weekly for
ninety minutes. Interruptions extended the nominal eight-week schedule to
approximately twelve weeks. Both conditions shared the first seven sessions,
developing reasons, elaborating them, using evidence, evaluating reasons,
and assembling an argument. Comparison activity ended after this phase;
it did not receive nine matched-duration alternative sessions. The full
condition continued with opposing-side reasons, counterarguments, rebuttals,
mixed evidence, practice dialogues and a final debate. The difference is
therefore a scaffolded package with more exposure, not a dialogue-only
manipulation at equal dose.
Comparison pupils also engaged in same-side team activity; this was not
an individual-only or no-collaboration control.

Adult coaches supported the activities. Participants made reason cards,
considered how an opponent would react, attached evidence and rebuttal cards,
and revised or dropped reasons. Some evidence articles were initially
presented without later opposing-perspective paragraphs, which were restored
in the mixed-evidence phase. Coaching, external representation, evidence
access, competitive goals and peer response are substantive ingredients.
The final debate used rotating speakers and short team consultations.
None was independently varied to identify one unique active component
\cite[pp.~1250--1251]{KuhnUdell2003Development}.

Individual assessments before and after activity elicited the participant's
opinion and supporting reasons, then all reasons that someone might offer
on either side. Unique reasons indexed topic-specific knowledge. Product
quality used a hierarchy: unsupported or weakly justificatory reasons,
conditions/administrative objections, reasons addressing the institution's
functions, and functional reasons comparing it with alternatives. An
argument could also become two-sided by including reasons on both sides.
That is the authors' evaluative rubric; an administrative injustice does
not become morally unimportant simply because the rubric places it below
a comparative functional reason. Seventy percent of product responses were
double coded: 81 percent agreement, Cohen \(\kappa=.67\).

Discourse assessments paired students with opposing opinions, usually with
the same partner at pre/post assessment, for videotaped ten-minute dialogues.
Utterances were segmented and coded by their relation to the opponent's
preceding move. \emph{Counter-A} introduced an alternative argument;
\emph{Counter-C} criticized that move. \emph{Clarify} elaborated one's
own position; \emph{Add} elaborated conversationally without advancing
the argumentative goal. A rebuttal sequence was a counter followed directly
by Counter-C. The full source appendix distinguishes further questions,
agreement, interpretation and other functions. Off-topic and meta-dialogue
utterances were excluded from the proportion analyses. Training used
15 percent of dialogues, and a fourth rater recoded 25 percent of the
remaining coded material: 90 percent agreement, \(\kappa=.72\).
Coding reliability concerns the specified categories, not truth of the
participants' substantive claims
\cite[pp.~1248--1249,1252--1253,1259--1260]{KuhnUdell2003Development}.

For the statistics in these records, \(F\) is a reported analysis-of-variance
statistic, \(t\) a reported t-test statistic and \(p\) a statistical p-value.
One experimental dyad's recording was lost, leaving 32 participants for
discourse analysis, nineteen experimental and thirteen comparison. Mean
utterance counts per participant/dialogue changed from 7.23, SD~5.06,
to 29.58, SD~17.56, in the full condition, versus 7.15, SD~4.25,
to 11.11, SD~9.83, in comparison. These are within-dialogue utterances,
not independent participant recruitments. To separate composition from
volume, the authors analyzed utterance proportions with an arcsin
transformation and pretest covariates.

The overall group effect was \(F(5,25)=6.0,p=.001\). Counter-C
proportions rose from 5.3 to 30.6 percent in the full condition and
11.4 to 21.4 percent in comparison; its adjusted group test was
\(F(1,29)=4.21,p=.049\). Counter-A changed from 7.6 to 13.6
versus 5.3 to 1.2 percent, \(F(1,29)=21.07,p=.001\).
Clarify declined from 37.2 to 18.5 versus 44.3 to 32.4 percent.
Simple disagreement did not meet .05 in its group test, \(p=.08\).
The source's effects are therefore not uniform across discourse categories
\cite[p.~1252, tables~3--4]{KuhnUdell2003Development}.

At least one Counter-C occurred for all nineteen experimental discourse
participants at posttest, versus ten of thirteen comparison participants;
initial counts were eight and five. Rebuttal-sequence means rose from .16,
SD~.50, to 3.79, SD~1.75, in the full condition, versus .23,
SD~.44, to .46, SD~.66, in comparison. The time/condition interaction
was \(F(1,30)=32.37,p<.001\). The paper prints no confidence intervals
for these contrasts and reports no dyad/team/school-adjusted outcome model.
The authors also stress that advanced strategies remained uncommon:
clarification requests and more complex sequences did not show a broad
uniform advance. Observed Counter-C usage is not mastery of every
argumentative skill
\cite[pp.~1253,1256--1257, table~6]{KuhnUdell2003Development}.

Individual-product quality used all 34 completers: positive change was
reported for 15/21 full-condition participants and 5/13 comparison
participants. The article describes that between-condition difference as
not reaching statistical significance. Eight of the fifteen experimental
improvements included a shift to two-sided argument. The knowledge-of-reasons
analysis instead omitted one experimental participant with a missing
all-reasons response. Per-assessment reason counts rose from 3.30 to 4.45
in that group and 2.31 to 2.54 in comparison, but the between-group gain
difference was not significant. Distinct reasons pooled across the two
assessments averaged 6.70 versus 4.38, with \(t(31)=3.23,p=.003\).
That pooled repertoire statistic is not the posttest gain or a new sample
\cite[pp.~1254--1255, table~7]{KuhnUdell2003Development}.

The authors themselves retain the possibility that extending Phase~I
could improve the comparison group because its total exposure was shorter.
This limits a claim that interaction with an opposing side alone caused
the difference. The outcomes concern one practiced topic and its pre/post
dialogue/product tasks. No delayed novel-topic transfer test is reported
in this experiment; transfers cited in its discussion belong to other
studies. The experiment supports changes under the implemented package,
with attrition, small samples and dependence retained, not a universal
age transition or a demonstration that every debate improves rationality
\cite[pp.~1256--1258]{KuhnUdell2003Development}.

\subsection{Kuhn and Crowell's 2011 study: extended practice and novel essays}
\label{sec:p19-r151-kuhn-crowell}

The 2011 intervention tests a different transfer: from sustained dialogic
activity to individually written essays on assessment topics not taught in
that activity. Entering sixth graders were eleven or twelve and attended
one academically challenging urban public middle school in the northeastern
United States. The neighborhood was ethnically diverse and low-to-middle
income; the article reports 80 percent Hispanic or African-American and
60 percent free/reduced-lunch eligibility. Individual age ranges beyond
the entry description, native-language eligibility and a probability
recruitment frame are not supplied.

The school formed three classes of thirty to thirty-two students matched
on gender, ethnicity and standardized scores. Two classes were randomly
selected for intervention and one for comparison. This is random assignment
at the level of only three classes, not independent randomization of
each student's curriculum. Table~\ref{tab:p19-r151-kuhn-crowell-samples}
keeps the overlapping and distinct outcome datasets separate
\cite[pp.~546--547]{KuhnCrowell2011Dialogic}.

\begin{table}[htbp]
\centering\small
\begin{tabularx}{\textwidth}{@{}Xrr@{}}
\hline
Dataset & Intervention & Comparison\\
\hline
Main longitudinal sample & 48 & 23\\
Expanded main Year~3 posttest sample & 60 & 28\\
External Year~3 comparison school & --- & 50\\
Following-entry replication & 37 & 21\\
\hline
\end{tabularx}
\caption{Student counts for distinct outcome samples in Kuhn--Crowell 2011. The first two
rows overlap and must not be added as independent recruited populations;
the external group was matched rather than randomly assigned.}
\label{tab:p19-r151-kuhn-crowell-samples}
\end{table}

In the longitudinal sample there were 27 girls/21 boys in intervention and
11 girls/12 boys in comparison. Most supplied at least three of four
assessments; one comparison student supplied only initial/final assessments.
Three missed the final assessment and seven missed one interim assessment.
The article does not specify an outcome-level missing-data/imputation
procedure that reproduces all printed inferential denominators. The
expanded posttest samples included pupils absent initially or arriving
later, with sex counts 31/29 and 13/15. The external fifty eighth graders
came from another city school matched on scores, lunch eligibility and
ethnic composition, not from the intervention's randomization.

Both curricula met twice weekly for fifty-minute philosophy classes.
The intervention spent about thirteen sessions per topic, four topics
per year, across three years. Same-side preparation generated and ranked
reasons. Partners then conducted electronic dialogues with six opposing
pairs, preserving chat transcripts for reflection. Group preparations,
a rotating-speaker debate, argument-map debriefing and a final essay
completed a topic cycle. Evidence materials were added from the final
topic of Year~1; students increasingly generated questions themselves,
with coaches supplying concise answers. Reflection changed form in
Year~3, and topics became three-sided. Note~2 records differences
between the two classes' Year~1 time allocation. This is a multicomponent,
changing curriculum rather than one constant dialogue-dose treatment.

Comparison students received teacher-led whole-class discussion, some
activities such as dramatization and more frequent essays: fourteen per
year versus four in intervention. Thus the control was educationally
active and had more writing practice, while coaching, peer structure,
reflection and evidence access differed. It does not isolate the causal
contribution of electronic rather than spoken dialogue or any single
ingredient
\cite[pp.~547--548,552, note~2]{KuhnCrowell2011Dialogic}.

Teacher-pay essays were assessed initially and after each of three years.
A separate assisted-dying essay appeared only at the Year~3 endpoint.
These were assessment subjects, not the curriculum's deeply studied
topics; neither political nor medical position was the answer key.
At Year~3, essays received twenty minutes and follow-up questions ten
minutes, with the teacher-pay assessment first on a different day.
The follow-up asked which unanswered questions would help the argument;
those assessment questions were not answered. More questions therefore
measure proposed evidence seeking, not actual verified information use.

The scoring categories preserve the source's distinctions. An own-side
argument offers advantages of the favored position. A \emph{dual perspective}
also attacks the opposing option, so all considerations can still favor
the same conclusion. An \emph{integrative perspective} acknowledges a
disadvantage of the favored position or advantage of its opponent, requiring
considerations leading in opposite directions to be addressed. This is
not identical to the 2003 paper's one-/two-sided argument coding.
Idea-unit coding was blind to condition and time; a random third of
essays was double coded. Agreement was 88 percent, \(\kappa=.76\),
for teacher pay and 93 percent, \(\kappa=.91\), for assisted dying.
Reliability of classification does not itself validate every factual
premise or normative recommendation within an essay
\cite[pp.~546,548--550, table~1]{KuhnCrowell2011Dialogic}.

For teacher-pay dual-perspective counts, the reported condition/time
interaction was \(F(3,67)=6.11,p<.001,\eta_p^2=.216\).
At Year~3, the percentages producing at least one dual-perspective
argument were 79 versus 29; integrative arguments appeared in 30 versus
zero percent, having been absent at earlier main-study assessments.
Both groups' essays became longer. Year~3 total argument means were
3.96 versus 2.61; the source also reports dual-perspective shares .49
versus .24 after accounting for overall production. Figure~1 error bars
are standard errors, not 95 percent confidence intervals.

On the new assisted-dying topic, dual-perspective argument counts averaged
1.17 versus .36, \(t(86)=3.85,p<.001\). Percentages supplying at
least one dual-perspective argument were 73, 29 and 24 for intervention,
internal comparison and external comparison; integrative percentages were
30, 18 and eight. These are different endpoints with different samples,
not universal success percentages. Proposed question counts averaged
4.18 versus .33 for teacher pay and 3.26 versus 1.64 for assisted dying,
with .92 in the external comparison. Counts and general-versus-case-specific
question coding support a limited evidence-seeking interpretation; the
article does not establish the future truth or usefulness of every question
\cite[pp.~549--551, tables~2--5]{KuhnCrowell2011Dialogic}.

The next incoming cohort provided a within-school replication through
Year~2, not an independent multi-population replication. After classroom
reassignment for behavior-management reasons, all students switching
conditions were excluded; note~4 does not give that exclusion count.
The retained sample was 37 intervention and 21 comparison students,
with comparison assessments available only at Year~2. Two intervention
responses were missing initially and two at Year~1. At Year~2,
dual-perspective count means were 1.05 versus .43,
\(t(56)=2.81,p=.007\), and integrative arguments appeared in 14
versus zero percent. Proposed question means were 1.48 versus .90,
but their difference was not significant, \(p=.148\). Replication
therefore does not mean every main-study outcome repeated at the same
strength or timing
\cite[pp.~550--552, table~6, note~4]{KuhnCrowell2011Dialogic}.

The paper supplies no interval estimates for these principal contrasts
and no reported class-cluster-adjusted model. Treating every pupil as
an independently randomized intervention unit would overstate the
randomization evidence. Attrition, changing membership, excluded
class-switchers, differences in adult facilitation and a restricted
school setting qualify causal and generalization claims. Nevertheless,
novel-topic individual essays provide an actual transfer outcome absent
from the undergraduate selection study and the same-topic 2003 experiment.
They do not establish transfer to abstract deduction, measured working-memory
capacity, an explicit concept of logical validity, all populations or
long-term performance after the curriculum ends.

\subsection{What these comparisons leave distinct}
\label{sec:p19-r151-social-capacity-boundary}

A complete selection under a stipulated conditional rule, a partner-responsive
counterargument, a higher-rated individual argument and an integrative essay
are different achievements. The first permits an exact truth-preservation
benchmark; the others use declared task-specific functional or evaluative
rubrics. The studies supply evidence at their measured times and scopes,
not one shared rationality score. The capacity and metalogical accounts in
this chapter additionally concern resource limits and explicit understanding.
These social-intervention outcomes do not directly measure either construct
or isolate capacity growth as their causal mechanism. A developmental
comparison must preserve what was practiced, what was tested, who was
sampled and whether a new topic or later occasion was actually observed.



\subsection{An original validity study: three tasks, three cohorts}
\label{sec:p19-r151-validity-record}

Moshman and Franks's 1986 paper investigated whether students distinguish the
necessary premise--conclusion relation from the empirical truth of its terms.
Experiment~1 recruited volunteers from Lincoln, Nebraska public schools:
20 fourth graders, ages 9 years to 10 years 6 months, mean 9 years 8 months,
and 20 seventh graders, ages 12 years to 14 years 3 months, mean 13 years.
The 41 University of Nebraska--Lincoln undergraduates came from a course
required of education majors; ages 18--37, mean 21.3 years. Undergraduates
reporting a logic course were not included; the count excluded is not given.
These are educational convenience samples, not representative age populations
\cite[pp.~153--156]{MoshmanFranks1986Validity}.

Seven arguments varied premise truth, conclusion truth and validity. Most used
transitive relations; a disjunctive inference prevented identification of all
valid items by identical form. Participants sorted argument trios and the whole
set and ranked the arguments. Children were individually interviewed with
movable cards, reading checks, follow-up questions and suggested sortings;
college students completed written booklets in small groups. That accommodation
reduces some expression demands but also makes the age comparison a comparison
of different administration modes. Validity coding required an explanation
distinguishing necessary consequence from truth, content or mere form; ranking
the three valid arguments above the others needed no verbal explanation.
Second coding of 50 written responses gave 82 percent category agreement and
98 percent agreement on the validity distinction.

The proportions using validity at least once were 0, 45 and 85 percent in the
fourth-grade, seventh-grade and undergraduate cohorts. The paper reports
\(\chi^2(2,N=81)=40.4,p<.001\). Zero successes among these 20 fourth
graders is an observation under the scoring and procedure, not proof that no
child of that age can possess a concept of validity. Among seventh graders,
individual responses ranged from fourth-grade-like to adult-like performance
\cite[pp.~156--158]{MoshmanFranks1986Validity}.

Experiment~2 used new volunteers: 18 fourth graders, ages 9 years 1 month to
11 years 4 months, mean 9 years 9 months; 18 seventh graders, ages 12 years
7 months to 13 years 11 months, mean 13 years; and 20 undergraduates, ages
19--23, mean 20.9. Recruitment used the same school and education-course
settings. Definitions, contrasting examples and explicit instructions preceded
40 validity judgments using five simple argument types with true, false or
neutral content. A separate ten-item control asked participants to choose the
better conclusion, rather than judge the validity of an entire argument.
The stringent criterion, at least 90 percent correct over 40 judgments, was
met by 11, 67 and 80 percent, respectively; a 70 percent criterion was met by
56, 89 and 100 percent. The control means were 9.6, 9.7 and 10.0 out of ten,
and nobody scored below eight. On the twelve critical items where truth and
validity conflicted, 11, 61 and 75 percent met the paper's above-chance validity
criterion of eleven or twelve correct
\cite[pp.~158--161]{MoshmanFranks1986Validity}.

Experiment~3 again used new volunteers: 20 fourth graders, mean 9 years
7 months, range 9 years to 10 years 5 months; 20 seventh graders, mean
13 years, range 12 years 5 months to 13 years 5 months; and 20 undergraduates,
mean 22 years, range 19--43. Within each age group, half were randomly assigned
an initial explanation of validity. All individually sorted the same forty
arguments with trial-by-trial correctness feedback and access to earlier cards.
Its success criterion was 90 percent correct over all forty, the last twenty,
\emph{or} the last ten judgments. That more permissive disjunctive rule is not
the identical measure used in Experiment~2. Success proportions were 10, 75
and 90 percent; \(\chi^2(2,N=60)=30.6,p<.001\). The paper does not report
confidence intervals for these proportions
\cite[pp.~161--164]{MoshmanFranks1986Validity}.

The authors favor a deficit in explicit validity understanding for most of
their fourth graders and a spontaneous-application limitation for many seventh
graders. They explicitly state that absent competence cannot be conclusively
proved and that new conditions might reveal it. Their evidence concerns several
cross-sectional cohorts, not an observed transition in each child between ages
ten and twelve. The strong control-task performance is decisive for keeping
ordinary inference separate from explicit validity classification. Nonsignificant
seventh-grade--college comparisons in the prompted tasks do not prove equal
population competence. The paper's internal contrasts test alternative task
explanations; they are not independent cross-cultural replications.

\section{Abstraction and biological proposals}
\label{sec:p19-r151-biological}

The chapter's Piagetian reconstruction links assimilation and accommodation
with reflective abstraction: common structure in lower-level activity becomes
available for more abstract operations. Its three broad levels are sensorimotor
action/perception schemas, concrete operations tied to objects, and formal
operations manipulating more general concepts. ``Locally rational'' means
coherence relative to the structures available at a level, not global conformity
to every logical norm. Markovits gives two opposing empirical pressures on a
simple stage interpretation: early successful deductions reported by Hawkins
and colleagues, and persistent adult errors. Overton and Ricco's
competence--performance distinction responds to the latter by separating
available logical capacity from conditions governing its use. The chapter
cites age-related instruction effects, but those effects do not independently
identify an unobserved competence structure
\cite[pp.~781--782]{Markovits2021DevelopmentRationality}.

Braine's natural-logic account proposes syntactic inference schemas, with a
biological argument that some rapid reliable inferences could serve survival.
Its interpretational component handles variability associated with pragmatics
and knowledge. This is neither a claim that every utterance is already in
formal notation nor a theorem that evolution installed a complete calculus.
Compare the full architecture of representation, schemas, control and pragmatic
integration in \cref{sec:p19-rationality-3-2}.

A different biological proposal attributes successful social-contract checking
to specialized cheater detection. The abstract selection task has four visible
cases: antecedent, nonantecedent, consequent and nonconsequent. For a stipulated
exceptionless conditional, the antecedent case can conceal a violation on its
reverse, and the nonconsequent case can conceal an antecedent. The remaining
two cases cannot conceal that particular violation. Thus checking a vowel--even
rule requires checking the vowel and odd-number cases. The chapter contrasts
poor abstract performance with improved performance under certain permission
or exchange rules and cites Cummins and Girotto and colleagues for children's
deontic versions. Agreement with the formal answer in that context does not
demonstrate domain-general logical competence or uniquely distinguish an evolved
module from learned pragmatic schemas. The drinking-age illustration is a
stipulated historical task rule, not a statement of current law
\cite[pp.~782--783]{Markovits2021DevelopmentRationality}.

Attachment supplies a further, explicitly different illustration of contextual
inference: expectations formed through interaction can organize later
interpretation of important others' behavior. Markovits connects Bowlby's
behavioral system, internal working models and adult romantic-attachment
research. These proposed social schemas do not establish truth-preserving
deduction; nor does a behavioral classification establish a particular caregiver
cause, a child's fixed future relationships, or an individual diagnosis.
The chapter's simplified caregiver examples must not be read as deterministic
causal laws or parenting advice. They illustrate experience-sensitive expectations,
not a measured developmental age trajectory
\cite[pp.~783--784]{Markovits2021DevelopmentRationality}.

\section{Conditional tasks and the capacity--content contrast}
\label{sec:p19-r151-conditionals}

\begin{hbconvention}{def:p19-r151-conditional-benchmark}{Necessity, invitation and uncertainty}
For the schematic discussion, \(P\) and \(Q\) are antecedent and consequent
propositions; \(A\) names an alternative antecedent. In an exceptionless
deductive reading of ``if \(P\), then \(Q\)'', asserting \(P\) supports
\(Q\) (MP), and denying \(Q\) supports denial of \(P\) (MT).
Asserting \(Q\) does not establish \(P\) (AC), and denying \(P\)
does not establish denial of \(Q\) (DA). A case with \(P\) false and
\(Q\) true is consistent with the conditional and defeats both invited
conclusions. Rejecting \emph{necessity} is not asserting that the invited
conclusion is false. A probabilistic or defeasible causal reading changes
the benchmark and must be declared separately.
\end{hbconvention}

Barrouillet's developmental mental-model explanation, as reconstructed by the
chapter, emphasizes the number of models working memory can maintain. The
initial represented case has both terms true. A second case with both false
gives the two possibilities of a biconditional; the additional case with a false
antecedent and true consequent supplies the full conditional possibilities.
Searching only the initial case favors MP and AC; the biconditional set favors
all four conclusions; the three-case set supports MP/MT but defeats AC/DA.
These are predictions of an assumed representation/search procedure, not
observed mental objects uniquely recovered from a response pattern. The
account also permits pragmatic interpretations that are appropriate to a
context without matching the intended formal benchmark; its strongest
capacity-based sequence must be distinguished from that qualification.

In the truth-table task, a participant classifies the four antecedent/consequent
combinations as making a rule true, false, or neither. This differs from
deciding whether an inference is necessary. Markovits reports a progression
from predominantly conjunctive judgments through biconditional judgments to
conditional judgments, associated with working-memory measures in selected
studies. The approximate ages he reports---eight or nine, eleven or twelve,
and mid-adolescence---describe those tasks and theoretical interpretations.
They do not license a universal calendar independent of language, content,
response categories or sampling. The chapter expressly cites the 2016
content-sensitive counterevidence
\cite[pp.~784--785]{Markovits2021DevelopmentRationality}.

\subsection{The original 2016 within-pupil comparison}
\label{sec:p19-r151-2016-record}

Markovits, Brisson and de Chantal tested 77 native French-speaking pupils
from an elementary school in the Montreal region: 35 in grade three,
20 girls and 15 boys, mean age 9 years 2 months; 42 in grade six,
21 girls and 21 boys, mean age 12 years 1 month. Entire classes received
booklets and common instructions. Each pupil evaluated two arbitrary
conditionals and two imaginary category--property conditionals; two booklet
versions reversed their order. Specific cases represented each of the four
antecedent/consequent combinations, with order varied. The paper does not
specify a probability recruitment frame, response rate, or interval estimates
for its pattern percentages
\cite[pp.~1908--1910]{MarkovitsBrissonDeChantal2016Conditionals}.

Arbitrary rules associated colors or clothing without an explicit relational
structure. The imaginary rules concerned category membership and properties
on an unfamiliar planet. The latter remove familiar factual knowledge while
retaining an interpretable category--property relation. This is not simply
``familiar versus unfamiliar words'': the rule types also differ in relational
structure and narrative framing. Responses were true, false or one-cannot-know.
A \emph{defective conditional} pattern marked the both-true case true, the
true-antecedent/false-consequent case false, and both false-antecedent cases
irrelevant. That pattern is not a complete material-conditional truth table,
which assigns true to its false-antecedent cases.

For grades three and six, respectively, defective-conditional pattern
percentages were 25.7 and 26.2 with imaginary rules versus 8.6 and 7.1 with
arbitrary rules. Conjunctive pattern percentages were 11.4 and 25.0 versus
22.6 and 39.3. Each pupil contributed two rule patterns per type; the reported
chi-square comparisons, \(\chi^2(1,N=70)=7.24,p<.01\) and
\(\chi^2(1,N=84)=8.85,p<.01\), therefore do not describe 70 and 84
independently recruited children. A within-subjects analysis of false responses
on the two false-antecedent combinations reported means 2.39 for arbitrary
and 1.71 for imaginary rules, with \(F(1,73)=14.08,p<.001\) for rule
type. Its degrees of freedom, and those in the separate basic-combination
analysis, do not provide a reported exclusion account; none is invented here
\cite[pp.~1910--1911, tables~1--2]{MarkovitsBrissonDeChantal2016Conditionals}.

The authors argue against a capacity-only explanation that preadolescents
cannot produce a more complex interpretation. They examine pragmatic-heuristic
and systematic-misinterpretation alternatives. Their contrast shows that a
uniform conjunctive limitation does not describe all these responses. It does
not directly measure working-memory capacity, prove that working memory never
matters, or establish identical comprehension in all children. High proportions
of incompatible and other patterns remain. The arbitrary-rule results partly
reproduce the earlier pattern, while the structured rules challenge its
unqualified generalization. This is a content-bound conceptual comparison,
not a longitudinal study or an independent multilingual replication
\cite[pp.~1911--1912]{MarkovitsBrissonDeChantal2016Conditionals}.

\section{Semantic retrieval, redescription and divergent thinking}
\label{sec:p19-r151-semantic}

Semantic redescription theory gives a different weight to knowledge. Premise
interpretation cues information in semantic memory, which can contribute
additional models. An alternative cause---another hard object breaking a
window---supports uncertainty about whether a rock caused an observed break.
A disabler---a resistant window---instead challenges the applicability of
the assumed causal rule. Retrieval is not necessarily restricted to the
information useful for the current inference: accessing alternatives can also
activate disablers. Thus improving search can lower acceptance of AC/DA yet
also lower acceptance of MP/MT. In a stipulated deductive task, inhibiting
premise-inconsistent information is appropriate to its benchmark; in an
empirical causal task, reconsidering the rule may be appropriate. The resulting
response cannot be evaluated without that task distinction
\cite[pp.~785--786]{Markovits2021DevelopmentRationality}; compare
\cref{sec:p19-rationality-5-4}.

The chapter distinguishes category--property, ordinary causal,
contrary-to-fact causal, and abstract conditionals. Its example relations are
dogs and legs, rocks and breaking windows, feathers and breaking windows,
and the uninterpreted terms XY and blrp. These vary the available meaningful
relation, not merely sentence length. Category knowledge can
supply an alternative animal with the named property; ordinary causal memory
can supply a different cause. Contrary-to-fact and uninterpretable abstract
rules require construction beyond that stored relation. Its reported pattern
places success with some category premises as early as seven or eight,
ordinary causal success later, contrary-to-fact success chiefly in later
adolescence, and consistently abstract success as difficult even for educated
adults. These are qualified summaries of selected studies, not required
success ages for each child. The theory proposes increasingly abstract
redescription of alternatives, rather than maturation alone. Markovits also
reports selective benefits of generating alternatives at a higher abstraction
level and interference from initially reasoning with abstract premises.
Those claims concern specified content and sequence, not a general instruction
that every additional example improves every reasoning task
\cite[pp.~785--786]{Markovits2021DevelopmentRationality}.

The divergent-thinking extension concerns generating possibilities beyond the
strong associations suggested by a premise. More ideas, more different
categories of ideas, and more original ideas are different measures. None is
identical to logical validity or to creativity in every domain. The chapter's
preschool example contrasts accepting a category inference with withholding
certainty about its converse, and discusses a preceding possibility-generation
exercise. The direct original records below delimit that claim.

\subsection{The original 2017 preschool records}
\label{sec:p19-r151-2017-record}

De Chantal and Markovits's Study~1 recruited 32 native French-speaking
children from a Montreal-region day-care center, 21 boys and 11 girls,
ages 41--64 months, mean 53.5 months. Parents consented. Individually
administered sessions in a quiet room lasted about fifteen minutes. Each child
completed idea generation, a dimensional-change card sort (DCCS), and reasoning,
with the first two tasks in one of two orders and reasoning always last.
Generation asked for ways of making noise, giving gifts and getting dirty,
with two examples for each prompt and no time limit. Flexibility was the mean
number of distinct response categories; fluency was mean number of responses.
Two category coders gave weighted kappa .798. The three-phase DCCS scored
threshold success in sorting by color, switching to shape, and switching under
a border rule. Although nominally zero to three, observed scores were only
one or two: all succeeded in phase one, none in phase three
\cite[pp.~210--214]{DeChantalMarkovits2017Alternatives}.

The reasoning task used three familiar universal category--property premises,
with AC followed by MP and requests for justifications. A point required both
accepting the necessary MP conclusion and withholding the AC conclusion for
the same premise; the total ranged from zero to three. Mean total was .81,
SD .90, equivalent to 27 percent of premise pairs, \emph{not} 27 percent
of children passing a universal competence criterion. Mean MP and AC scores
were 2.91, SD .39, and .91, SD .89, respectively. The published correlations
with total reasoning were .48 for flexibility, .37 for fluency and -.06 for
DCCS. Controlling fluency, the flexibility correlation was
\(r(29)=.353,p=.05\). These are small-sample associations, not intervention
effects; the restricted DCCS score range and contested construct interpretation
limit an inference that inhibition has no role. No confidence intervals for
these correlations are printed
\cite[pp.~212--214, tables~1--2]{DeChantalMarkovits2017Alternatives}.

Study~2 recruited another 32 native French-speaking children from a
Montreal-region day-care center, 16 boys and 16 girls, ages 32--61 months,
mean 47.06 months. Children were randomly assigned generation before reasoning
or the standard two-phase DCCS before reasoning. The latter was shortened
to balance administration length and because no Study~1 child succeeded in
phase three. Five category--property premise pairs replaced the three-item
task; cactus content was removed, and new content was added. The mean paired
reasoning scores, out of five, were 1.50, SD 1.93, after generation and .38,
SD .81 in table~3, after DCCS; the prose prints .80 for the latter SD.
An age-adjusted analysis gave \(F(1,29)=4.39,p<.05\), partial
\(\eta^2=.132\). The study reports no confidence interval for that contrast
\cite[pp.~214--216]{DeChantalMarkovits2017Alternatives}.

This random assignment supports a comparison between those two preceding
activities in this brief task. It lacks a neutral no-prime condition and a
pre/post reasoning measure, so it cannot separately identify facilitation
from generation and impairment from DCCS. It does not measure lasting transfer
or an individual developmental transition. Near-ceiling MP acceptance leaves
most score variation in AC uncertainty. The authors discuss a yes-response
bias as potentially inflating MP success and depressing AC success; they
explicitly decline to conclude that inhibition is unimportant. Their
justification analyses count multiple responses per child, not independent
recruitments. The English boots-question gloss on p.215 repeats the
consequent rather than asking about boots; the French appendix on p.218 asks
whether the footwear is boots. The French protocol is the reported administered
language, and the published scores are retained rather than recalculated from
that mistranslated gloss
\cite[pp.~215--218]{DeChantalMarkovits2017Alternatives}.

The 2020 primary is titled \emph{Divergent Thinking Promotes Deductive
Reasoning in Preschoolers}, rather than the title used in the chapter's
reference list. Its author-reproduced main record follows below with its
own population, protocol and version qualifications. The separately cited
2012 school-age study compares socioeconomic groups; that original study
is not supplied by the preschool records. The chapter-level illustration of
generation without examples must not be retroactively assigned to the 2017
procedure, which supplied examples. Estimates for one population or procedure
cannot be borrowed to complete the other study's record
\cite[pp.~786--789]{Markovits2021DevelopmentRationality}
\cite{DeChantalGagnonMarkovits2020PreschoolMetadata}.

% Reader-ready C05 evidence continuation for ED26-BM-RAT2021-15-1.
% Root-owned integration: these records extend the semantic/divergent-thinking
% section; core chapter and original intakes were intentionally not edited.
\subsection{The preschool experiment published in 2020: an author-reproduced record}
\label{sec:p19-r151-2020-author-record}

The 2020 article by de Chantal, Gagnon-St-Pierre and Markovits is a different
study from the two 2017 cohorts. Its complete author-reproduced main text is
Chapter~VIII, pp.~95--134, of de Chantal's August 2021 UQAM thesis.
The thesis's own bibliography identifies the published article as
\emph{Divergent Thinking Promotes Deductive Reasoning in Preschoolers},
\emph{Child Development} 91(4), 1081--1097, DOI 10.1111/cdev.13278.
The following methods and results are attributed to that exact thesis
presentation, not asserted byte-identical to the unavailable journal PDF.
The reproduced text refers to an appendix of full reasoning stimuli and
category coding examples that is not included with this study in the thesis;
those missing protocol details are not borrowed from its separate 2017 appendix
\cite[pp.~92--134,151]{DeChantal2021ThesisSelectedStudy}.

The experiment tested whether a preceding generation exercise helps children
distinguish a necessary inference from an invited but uncertain converse,
and whether explicit examples alter that benefit. Its comparison was between
generation without examples (DT-free), generation with examples (DT-examples),
and a preference-report control. This addresses the lack of a neutral
comparison in the 2017 generation-versus-card-sort study; it does not create
a longitudinal developmental test.

One hundred twenty children were recruited through eight day-care centers
in the Montreal region. The reported age range was 46--79 months, mean
60.0 months, SD~5.94; there were 65 boys and 55 girls. That range extends
outside the abstract's nominal four-to-five-year description. The sample
was predominantly Caucasian and drawn from middle-class neighborhoods.
Parental consent confirmed that French was mainly spoken at home; three
children also spoke another language and were reported to complete the
experiment without difficulty. This is a restricted local recruitment frame,
not a representative sample of all preschool children
\cite[pp.~107--112]{DeChantal2021ThesisSelectedStudy}.

Four experimenters received common training and were blind to the hypothesis.
Individual sessions used an animated, recorded presentation in a separate
day-care room. The experimenter sat beside the child, wrote responses,
handled questions and restarted the video after each required answer. A
white board limited surrounding distractions and possible environmental cues.
Children were randomly assigned one of the three preliminary activities,
then received the reasoning task and finally a knowledge/comprehension check.
Sessions lasted seven to twelve minutes and ended with a sticker. Hypothesis
blinding does not establish condition blinding: the example manipulation was
part of the presented activity.

Generation prompts asked for ways of making noise, giving gifts and getting
dirty, always in that order. DT-examples supplied two examples for each
prompt: a whistle/paper tearing, a soft toy/chocolate, and mud/ketchup.
DT-free used the same themes but requests for more possibilities without
those examples. The control instead asked about favorite colors, games and
foods, also in fixed order, to elicit roughly comparable verbal output.
Experimenters could not give feedback or encouragement but could prompt for
different ideas after a pause of approximately five seconds, and moved on
only when the child reported having no more ideas. An exercise without
examples was therefore not an absence of instruction or interaction
\cite[pp.~108--111]{DeChantal2021ThesisSelectedStudy}.

The reasoning task had four familiar category--property premises, formulated
as universal statements. Each received MP first, then AC, with standardized
requests for explanations. This reverses the AC-before-MP sequence in the
2017 record. The main text's examples and justifications concern dogs/legs,
trucks/wheels, flies/wings and oranges/edibility; they do not supply the
missing appendix's complete wording. The subsequent check explicitly asked
for an alternative to each premise's category, stopping after one answer.
Passing such a prompted retrieval check does not demonstrate spontaneous
retrieval during the earlier inference task.

Three children had incomplete sessions, two departed from the protocol,
five were excluded following educators' reports of verbal production or
comprehension difficulties, and seven failed the four check questions.
The latter exclusions were four from DT-examples, one from DT-free and
two from control. The analyzed cohort was therefore 103, not 120:
51 boys and 52 girls, age 46--79 months, mean~60.3, SD~5.69.
Final group sizes were 32, 32 and 39 for DT-examples, DT-free and control,
respectively. Their mean ages were 59.3, 60.2 and 61.1 months, with SDs
5.60, 5.96 and 5.55. The main text does not give each initially randomized
group's size or a full condition breakdown of the other ten exclusions;
equal initial allocation is not invented. The authors report a similar
pattern when the seven check failures were included, but do not supply
the corresponding numerical estimates here
\cite[p.~112]{DeChantal2021ThesisSelectedStudy}.

In this record, \(p\) denotes a statistical p-value, \(F\) an analysis-of-variance
statistic and \(\chi^2\) a reported chi-square statistic, not a logical
antecedent, connective or representation model. A correct MP answer expressed
certainty; a correct AC answer expressly denied certainty. ``I don't know''
was scored incorrect for both. Each separate inference score ranged from
zero to four. The paired deductive score awarded a point only when both
answers for the same premise were correct, also ranging from zero to four.
Overall MP and AC means were 3.87, SD~.50, and 1.25, SD~1.61;
paired-score mean was 1.18, SD~1.56, with mode zero. Unlike the 2008
categorical task below, a nonspecific uncertainty response was not grouped
with rejection. Thus changing score rules can change apparent competence.

An age-adjusted Poisson loglinear analysis reported a condition effect,
Wald \(\chi^2(2)=13.63,p<.01\); its age term was
\(\chi^2(1)=2.79,p=.10\). Paired-score means were 1.97, SD~1.75,
for DT-free, 1.09, SD~1.47, for DT-examples, and .62, SD~1.21,
for control. Bonferroni comparisons gave DT-free versus control
\(p<.001\), DT-free versus examples \(p<.05\), and examples versus
control \(p=.06\). These are mean counts out of four premise pairs,
not percentages of children with a domain-general competence. Figure~8.1's
error bars are standard errors, not confidence intervals for those contrasts
\cite[pp.~112--114]{DeChantal2021ThesisSelectedStudy}.

The source adds a log-score ANCOVA, reporting
\(F(2,99)=7.99,p<.01,\eta_p^2=.139\). It does not specify how zeros
were handled before a logarithm; no offset or exclusion rule is supplied here.
Its age-term line literally prints \texttt{p = 12}, which is not a valid
p-value and is not silently repaired. A nonsignificant goodness-of-fit
comparison does not establish that a bounded zero-to-four outcome has an
exact Poisson law. These reporting and model boundaries limit exact
reproduction without invalidating the distinction among the observed group
means; the Poisson and ANCOVA procedures are analyses of the same data,
not two independent replications.

The generation scores also need their operational definitions. Redundant
ideas and literal repetitions of supplied examples were removed. Labels
were standardized to avoid counting paraphrases as different ideas. The
two DT conditions were pooled to construct one frequency lexicon per theme,
yielding 93, 133 and 90 distinct ideas. Research assistants formed thematic
categories, and two coders assigned ideas to them, resolving disagreements
after Cohen kappa values .77, .75 and .64. Flexibility counted represented
categories; each uncategorized idea received its own flexibility point.
Originality counted ideas occurring once in that pooled sample divided by
the child's fluency, then averaged over the three exercises. This is
sample-relative statistical rarity, not an invariant property of an idea.
The main text does not state the division convention for an exercise with
zero ideas or reproduce the promised category appendix
\cite[pp.~116--119]{DeChantal2021ThesisSelectedStudy}.

Table~8.2 prints fluency, flexibility and originality means of
3.13, 1.98 and .19 in DT-examples, versus 3.45, 2.05 and .30 in
DT-free; corresponding SDs are 2.21, 1.43 and .21 versus 1.72,
.81 and .17. The earlier verbal-output paragraph instead prints 3.12
and 3.47 for the DT fluency means, and 3.23 for control. Those small
reporting differences are retained rather than declared identical.
The three-condition verbal-output analysis had \(F(2,99)=.23,p=.80\);
nonsignificance is not proof of equal verbal activation. In the DT-only
age-adjusted analyses, originality differed,
\(F(1,61)=5.49,p<.05,\eta_p^2=.083\), whereas fluency and flexibility
comparisons had \(p=.21\) and \(p=.38\). Reported Cronbach alpha across
three exercises was .74 for fluency, .71 for flexibility and only .40 for
the originality ratio; the undivided unique-idea count had alpha .67.
The lower ratio reliability matters to individual-difference claims
\cite[pp.~116--120, table~8.2]{DeChantal2021ThesisSelectedStudy}.

Age-controlled correlations of reasoning with fluency, flexibility and
originality were .45, .48 and .36 in DT-examples, versus .15, .10 and
.02 in DT-free. The reported between-condition correlation comparisons
did not pass .05: their p-values were .10, .054 and .08, respectively.
Significance in one group but not another does not itself establish different
population correlations. The proposed attenuation of individual differences
by DT-free is an interpretation, not demonstrated causal mediation by
originality. No interval estimates for these correlations or the principal
group contrasts are printed
\cite[pp.~120--126, table~8.3]{DeChantal2021ThesisSelectedStudy}.

The explanation data provide a further but dependent observation. Alternative
antecedents occurred in 55.0 percent of explanations following correct AC
answers versus 7.6 percent following incorrect ones. Disablers occurred in
83.3 percent of explanations following the relatively rare incorrect MP
answers. The reported justification-analysis denominators, 392 and 408,
count responses from the 103 children, not new participant samples.
Explanations elicited after judgments do not uniquely identify the process
that produced the judgment, and the study was not designed to identify
condition-specific qualitative changes in those explanations
\cite[pp.~114--116,125]{DeChantal2021ThesisSelectedStudy}.

Randomized preliminary activities support a task-specific short-term
comparison after the stated exclusions. The stronger no-example result
constrains a mere-more-answers explanation, but does not isolate originality
from every difference in semantic cueing, attention, motivation or task frame.
The authors explicitly leave duration, broader individual differences and
socioeconomic generalization open. Nothing here establishes a stable age
equivalence, transfer to abstract or causal premises, enduring improvement,
or a universally effective educational practice. The absent original appendix
and journal-version comparison remain material limits on exact protocol
qualification, not reasons to substitute the abstract for this evidence.

\subsection{The 2008 contrast: necessity, likelihood and confidence}
\label{sec:p19-r151-2008-record}

Markovits and Thompson's three original studies separate three questions:
whether a conclusion is necessary, how likely it is, and how confident the
reasoner is in the likelihood judgment. Their categorical answers and
five-point ratings are different measures. The paper's proposed two-strategy
interpretation is a psychological hypothesis, not a consequence of having
two response formats
\cite[pp.~1066--1069]{MarkovitsThompson2008Development}.

Study~1 involved 56 children from rural state schools in southwest England:
17 in Level~1, mean age 6 years 5 months, 2 girls/15 boys;
13 in Level~2, mean 7 years 1 month, 5 girls/8 boys; and 26 in
Level~4, mean 8 years 7 months, 13 girls/13 boys. The paper calls these
nominal six-, seven- and nine-year cohorts; it does not report individual
age ranges, a probability recruitment frame or response rates. Four
videotaped sequences varied premise order and whether categorical or
probabilistic questions came first, systematically from child to child,
not by a stated random allocation. Individual sessions lasted about ten
minutes. MP always preceded AC for a given premise.

Premises concerned cars/motors and hedgehogs/four legs, chosen to provide
accessible alternatives and few disablers. Categorical instructions asked
for certainty assuming the premise true, followed by a justification.
``Yes'' counted as acceptance; both ``no'' and ``not sure'' counted as
nonacceptance. The likelihood task used five increasing red bars, labeled
from not sure to very, very sure. Practice displays contained one, five
or nine red trees out of ten, with repeat instruction when needed.
The article supplies English formulations; it does not separately state
native-language eligibility for these English-school cohorts
\cite[pp.~1069--1071]{MarkovitsThompson2008Development}.

\begin{table}[htbp]
\centering\small
\begin{tabular}{lrrrr}
\hline
Level & Pupils & MP accept (\%) & AC accept (\%) & AC rating\\
\hline
1 & 17 & 88.2 & 61.8 & 3.57\\
2 & 13 & 88.5 & 11.5 & 3.38\\
4 & 26 & 84.6 & 25.0 & 2.72\\
\hline
\end{tabular}
\caption{Study~1's categorical acceptance percentages and separate AC
likelihood ratings on a one-to-five scale, combined over two premises.
The likelihood numbers are not numerical probabilities or correctness scores.}
\label{tab:p19-r151-2008-study1}
\end{table}

Table~\ref{tab:p19-r151-2008-study1} shows that categorical AC acceptance
does not form a strictly monotone age ladder in these small cohorts.
Mean MP likelihood ratings were 4.56, 4.50 and 4.56, respectively.
An instruction-task comparison therefore need not imply that all of one
cohort's inferences are superior. Counterexample explanations were coded
by two experimenters with fewer than five percent initial disagreements,
resolved by discussion. Their correlations with categorical AC acceptance
were -.69, -.61 and -.75 across the three levels; with AC likelihood
ratings they were -.56, -.59 and -.01. The paper reports a within-oldest-
cohort difference between those correlations. These relationships involve
elicited explanations, not randomized counterexample use, and the paper
prints no confidence intervals for them
\cite[pp.~1071--1073, tables~1--3]{MarkovitsThompson2008Development}.

Study~2 used another 53 pupils from the same kind of southwest-English
schools: 21 in Level~1, mean 6 years 5 months, 11 girls/10 boys;
18 in Level~2, mean 7 years 7 months, 8 girls/10 boys; and 14 in
Level~4, mean 9 years 8 months, 8 girls/6 boys. The visible high-likelihood
array contained four red squares, one red circle and four blue circles;
the low-likelihood array contained one red triangle, four red circles and
four blue circles. The rules linked squares or triangles to red.
Conditional on a red shape and equal selection among those red shapes,
the target proportions are four-fifths and one-fifth. Both arrays still
contain a red counterexample to the converse, so neither AC conclusion
is deductively necessary. Actual selection probabilities require that
equal-selection assumption, not just counting the picture.

The procedure repeated the four systematically varied sequence formats
but omitted justification requests. Average AC likelihood ratings were
3.87 for the high-likelihood array versus 2.56 for the low-likelihood
array, \(F(1,47)=22.64,p<.001\). This supports discrimination, not
numerical calibration of a probability distribution. Categorical AC
acceptance averaged 49.1 versus 28.3 percent by array. It also differed
with task order: 25.9 percent when categorical questions came first,
versus 51.9 percent when likelihood questions came first. The reported
order comparison was \(\chi^2(1)=6.34,p<.01\). Likelihood ratings
showed no reported order effect; that nonsignificance does not prove exact
order invariance. Grade-specific order cells contained only seven to
eleven children, and the stimuli/instructions differ from the preschool
generation experiment
\cite[pp.~1073--1075, table~4]{MarkovitsThompson2008Development}.

The printed probabilistic example on p.~1074 repeats the triangle
antecedent in its question instead of asking about the consequent, unlike
the surrounding inference description. Source-image inspection confirms
this wording discrepancy; it is not repaired by inventing a verified
administration script. The record preserves the stated paradigm and
published results while leaving that exact example's protocol standing
qualified.

Study~3 tested a separate 26 native French-speaking Grade~2 children in
a Brussels public elementary school, mean age 7 years 3 months,
13 girls and 13 boys. Two drawings showed eight red squares and one
blue square, or eight circles and one triangle. Each child first rated
the likelihood of choosing the frequent kind with eyes closed, then
confidence in that rating; the rare-kind questions followed. The same
five-point scale was used for both judgments. Drawing order varied
between two halves of the cohort; frequency-question order did not.
Sessions lasted approximately five minutes.

For the square display, likelihood means for frequent/rare outcomes were
3.61/2.62, whereas corresponding confidence means were 3.81/3.92.
For circles/triangle, the values were 3.54/2.88 and 3.69/3.96.
A low-likelihood event can therefore receive a relatively confident
rating without being rated likely. The rare-event likelihood/confidence
comparisons were \(t(25)=3.45,p<.001\) and
\(t(25)=3.11,p<.01\). The main likelihood ANOVA prints
\(F(1,23)=7.11,p<.01\); its denominator is retained rather than
used to invent an unreported exclusion record. The article gives no
intervals for the displayed means. This is evidence for a distinction
within this instructed rating task, not a demonstration that every young
child distinguishes confidence from probability in every context
\cite[pp.~1075--1076, table~5]{MarkovitsThompson2008Development}.

The article itself qualifies a simple analytic/heuristic dichotomy:
retrieving an alternative can be automatic when strongly associated,
but effortful when weakly associated; explicit representation is another
process. Its cross-sectional task differences do not measure maturation
of one pure system. Comparing these records with preschool-generation
scores also requires different recruitment, premise content, uncertainty
coding and score denominators to remain visible. Neither the 2008
instruction contrast nor the later generation prime establishes an
age-stage equivalence or a universally optimal reasoning strategy
\cite[pp.~1076--1078]{MarkovitsThompson2008Development}.



\section{Using the developmental account}
\label{sec:p19-r151-use}

\begin{hbexample}{def:p19-r151-three-outcomes}{A correct deduction, uncertainty and explicit explanation}
Stipulate that every dog in a described group has legs. A named dog must have
legs under that stipulation. A different animal known only to have legs might
be a dog or a cat; a legged cat is compatible with the stated premise, so
withholding certainty that this animal is a dog is appropriate. A third question asks the
participant to explain why the first inference holds even if the stipulated
category relation is empirically false. That asks for an explicit
truth--validity distinction. Success on the first question, success on both
inference questions, and success on the explanatory question are three outcomes.
Changing the premise to an ordinary defeasible expectation creates a fourth
task; mentioning an exception now challenges its scope rather than necessarily
showing a deductive mistake. This is a stipulated contrast, not measured data.
\end{hbexample}

The mental-model account in \cref{sec:p19-rationality-2-3} supplies the broader
representation and counterexample-search theory. The natural-logic account in
\cref{sec:p19-rationality-3-2} supplies schemas plus interpretation and control.
Neither general account is this chapter's complete developmental evidence.
Bayesian reasoning in \cref{sec:p19-rationality-4-5} evaluates uncertain
inference under an identified probabilistic interpretation; the developmental
tasks here often evaluate necessity under stipulation. A lower rate of
deductive answers after retrieving real-world exceptions is not by itself
evidence against every probabilistic or context-sensitive account.

Likewise, monitoring in \cref{sec:p19-rationality-8-6} distinguishes judgments,
accuracy and control. The adult paired-noun studies in
\cref{sec:p19-judgments-of-learning} predict later recall, not children's
deductive development. An improvement in delayed memory judgments cannot
discharge the age, language, premise-interpretation or transfer questions here.

\begin{hbmethod}{def:p19-r151-lookup}{Interpret a developmental reasoning result}
Specify whether the outcome is an inference, a truth-table judgment, an
explanation, an argument strategy or a monitoring report. State the premise
reading and response criterion, recruitment and age distribution, language,
instruction and task order, response opportunities, exclusions and score.
Preserve participant and item dependence and distinguish observed age
differences, within-person content contrasts and randomized activity effects.
Report dispersion and uncertainty where available; leave missing intervals
missing. Compare at least one plausible task explanation: knowledge access,
premise interpretation, response bias, expression demands, motivation or
resource demands. A failed arbitrary task is not sufficient evidence for
absent domain-general competence; a successful familiar task is not sufficient
evidence for universal transfer. The evidence-record distinctions in
\cref{sec:p19-evidence-methods} govern these inferences.
\end{hbmethod}

Markovits's final argument is that increasing knowledge has two roles.
It provides the alternatives and structure needed for more powerful reasoning,
while also supplying expectations and commitments that can pull judgments away
from a stipulated logical benchmark. Increasing computational resources and
increasing understanding of the empirical world can therefore coexist with
both more capable deduction and more context-sensitive departures from it.
This interaction is the chapter's synthesis proposal; it does not resolve
which mechanism dominates each task or supply a single rationality scale
\cite[p.~787]{Markovits2021DevelopmentRationality}.

\begin{proofboundary}{def:p19-r151-evidence-boundary}{What the evidence does and does not establish}
The complete 2021 chapter provides the attributed developmental comparison.
The detailed original papers supply specific validity, conditional
interpretation, preschool-generation and uncertainty-task records, with their
different samples and measures. The shared collaborative-selection record
and the two adolescent dialogic interventions distinguish immediate group
outputs, practiced discourse and novel-topic individual transfer, without
identifying these as one measured capacity. The 2020 thesis-reproduced study has its own
version and missing-appendix qualifications. Other cited findings remain chapter-level reports here, not
invented complete quantitative study records. No meta-analysis, participant-data
reanalysis, independent replication, longitudinal universal sequence or clinical
assessment is established. The logical case explanations fix benchmarks; no
mathematical completeness theorem for human cognition is promised. Neither a
task failure nor a favorable short-term activity comparison warrants individual
diagnosis, parenting advice, or a universally effective educational intervention.
\end{proofboundary}
